\section{The shared-point quartic in 962 variables}
\label{sec:v6-962}

This section gives a new complete quadratic compilation of the shared-point
construction in the public agentic notes~\cite{agentic-improvements}.
The shared comparison points, reciprocal-trace detectors, fixed height
scales, direct Lucas test, and affine anchor change are arithmetic
and algebraic improvements recorded in those notes. The simultaneous prime theorem used to choose
the shared points is Sun's Theorem~4.1, in the precise form of
Theorem~\ref{thm:v5-simultaneous-primes} below. Our contribution here is an
independently specified gate program, its source verification, and
further changes of coordinates and circuit identities leading to a
962-variable polynomial after 32 certified eliminations. The
definition does not depend on an unavailable coefficient table.
The relevant passages in the public notes are
\href{https://github.com/JoseBrox/diophantine-z-q_improvements/blob/f4d1b5e6c23d4adeb3ae1f986ab0bf87884bdff6/paper/diophantine_optimization_20261106_v3.tex#L2354-L3044}{shared-point construction},
\href{https://github.com/JoseBrox/diophantine-z-q_improvements/blob/f4d1b5e6c23d4adeb3ae1f986ab0bf87884bdff6/paper/diophantine_optimization_20261106_v3.tex#L3352-L3405}{fixed height scales},
and
\href{https://github.com/JoseBrox/diophantine-z-q_improvements/blob/f4d1b5e6c23d4adeb3ae1f986ab0bf87884bdff6/paper/diophantine_optimization_20261106_v3.tex#L4135-L4758}{quadratic compilation and anchor translation}.

We retain the arithmetic conclusions and conventions of
Section~\ref{sec:degree8-original}, including the restricted-converse
input to Lemma~\ref{hyp:o8-pole}. In this section the height estimate is
Theorem~\ref{hyp:o8-height}, with exponent 24 and its original constant
$H_*$; it is not the exponent-19 estimate used for the preceding quartic.
Thus
\[
  1\le H_*<2^{4783227}.
\]
The exact finite compilation and coefficient assertions below require no
arithmetic conjecture. Identifying the rational input projection with
$\Z$ uses the stated arithmetic results and the simultaneous prime theorem.

The point, pole, and correction variables below specify the
arithmetic source predicate. The final compilation subsection
gives explicit maps to its implementation coordinates.

\subsection{Points, local parameters, and quotient domains}
\label{subsec:v5-point-source}

Use the normalized curve and generator already established in
Section~\ref{sec:degree8-original}:
\[
 K=\Q(\sqrt2),\quad d=3-\sqrt2,\quad D=d^2,
 \quad E:y^2=x^3-Dx,
\]
\[
 P_*=(57/8-7\sqrt2/2,-165/8+417\sqrt2/32),
 \qquad 2E(K)=\Z P_*.
\]
The last assertion follows from Proposition~\ref{prop:generator} by the
unit change of variables recorded there. Conjugation of $K/\Q$ is
denoted by $\sigma$. All field equations below mean both rational
component equations in the basis $1,\sqrt2$.

The source uses nine nonidentity points
\[
 A,S,T,Z,H,U_A,U_S,U_T,U_Z\in2E(K).
\]
Each is represented as the double of an affine halfpoint
$(\alpha,\eta)\in E(K)$, using four rational halfpoint coordinates.
The equations for a double are
\begin{equation}\label{eq:v5-double}
 \eta^2=\alpha^3-D\alpha,\qquad
 2\eta m=3\alpha^2-D,\qquad
 x=m^2-2\alpha,\qquad y=m(\alpha-x)-\eta.
\end{equation}
Products of degree greater than two in this display are evaluated by
the quadratic gates specified below. Set
\[
 Q_0=P_*,\quad(Q_1,Q_2,Q_3,Q_4)=(A,S,T,Z),
 \quad(U_0,U_1,U_2,U_3,U_4)=(H,U_A,U_S,U_T,U_Z).
\]
Compute $[4]Q_i$ by two doublings, or by fixed rational arithmetic for
$Q_0$, and compute
\[
 V_i=U_i-[4]Q_i\quad(0\le i\le4).
\]
For the subtraction of $(x',y')$ from $(x,y)$ we retain an inverse of
$x-x'$ in $K$ and use
\[
 m=(y+y')/(x-x'),\quad x''=m^2-x-x',\quad
 y''=m(x-x'')-y.
\]

\begin{lemma}[Domains of the point quotients]\label{lem:v5-point-domains}
The equations \eqref{eq:v5-double} force $\eta\ne0$ and describe a
finite nonidentity point of $2E(K)$. Every such point has nonzero
abscissa and ordinate. The guarded differences $V_i$ have the same
properties. Consequently $f(Q)=1/x(Q)$ and $z(Q)=-x(Q)/y(Q)$ are
defined and nonzero at every point where they are used here.
\end{lemma}
\begin{proof}
If $\eta=0$, the curve equation gives $\alpha=0$ or $\alpha^2=D$.
The right side of the slope equation is then $-D$ or $2D$, respectively,
a contradiction. The ordinary duplication formulas therefore apply.
They give
\[
 x([2](\alpha,\eta))=
 \left(\frac{\alpha^2+D}{2\eta}\right)^2\ne0,
\]
because $D=d^2$ and the numerator is positive at either real embedding.
An ordinate-zero point has abscissa $0,d$, or $-d$.
The latter two are not squares in $K$: their rational norm is 7,
whereas the norm of a square is a rational square. Thus the doubled
ordinate is nonzero. A guarded affine difference is finite and is
again in $2E(K)$, so the same argument applies. Conversely any
nonidentity point of $2E(K)$ has an affine halfpoint of nonzero
ordinate and is represented by these formulas.
\end{proof}

Write
\[
 a_0=1,\quad a_1=-2/5,\quad
 a_j=\frac{4j-5}{4j+1}a_{j-1}\ (j\ge2),\qquad
 G_J(w)=\sum_{j=0}^J a_jw^j,
\]
and, for $e=16,52$, put
\begin{equation}\label{eq:v5-log}
 \ell_e(Q)=z(Q)G_{e/4-1}(Df(Q)^2).
\end{equation}
Lemma~\ref{lem:quartic-reciprocal-nonzero} applies to both truncation
degrees 3 and 12. Thus all these values are nonzero globally.

\begin{lemma}[The logarithm estimate at unramified places]
\label{lem:v5-unramified-log}
Let $p>76176$ be a rational prime and $\mathfrak p\mid p$ a place
of $K$. If a nonidentity point $Q$ is in the formal group at
$\mathfrak p$, normalize $v(p)=1$ and put $r=v(z(Q))$. Then $r\ge1$,
and for $e\in\{16,52,64\}$,
\[
 v(\log_E Q)=r,\qquad
 v(\log_E Q-\ell_e(Q))\ge(e+1)r,\qquad
 \ell_e(Q)/\log_E Q\in1+p^e\mathcal O_{K_{\mathfrak p}}.
\]
Here $\ell_{64}$ is defined by the same formula with $G_{15}$.
Ratios with numerator precision $e$ and denominator precision $e'$
in this set have the valuation of their full logarithm ratio and
relative error of valuation at least $\min(e,e')$.
\end{lemma}
\begin{proof}
The local extension is unramified, including when $p$ is inert.
The curve has good reduction and $D$ is a unit. On the formal branch,
\[
 f=z^2(1-Df^2),\qquad w=Df^2,\qquad v(w)=4r.
\]
The recurrence gives
$4w(1-w)G'(w)+(1+w)G(w)=1-w$ and hence the same identity
$\log_E Q=zG(w)$ as in Lemma~\ref{lem:quartic-reciprocal-log}.
For $j\ge1$ the denominator of $a_j$ divides $(4j+1)!$, so
\[
 v(a_j)\ge-\frac{4j+1}{p-1}.
\]
Set $m=e/4\le16$. If $j\ge m$ and $4j+1<p$, then
$v(a_jw^j)\ge4mr$. Otherwise $j\ge19044$ and
\[
 v(a_jw^j)\ge4jr-j\ge3jr\ge4mr.
\]
The lower bounds tend to infinity, proving the estimate for the
entire tail. Every nonconstant term of $G(w)$ has positive valuation,
so the leading term gives $v(\log_E Q)=r$. Division proves the
relative estimate, and division of two relative estimates proves
the ratio assertion. No step uses residue degree one.
\end{proof}

\subsection{The complete rational source}
\label{subsec:v5-complete-source}

The ordinary input is $X$, and $\beta=X^2+2$. Define
\[
 (k_1,k_2,k_3,k_4)=(16,52,52,52),\qquad
 L_i=\frac{\ell_{k_i}(U_i)}{\ell_{52}(H)},\qquad
 L'_i=\sigma\!\left(\frac{\ell_{k_i}(V_i)}{\ell_{52}(V_0)}\right).
\]
The sixteen field tests $D_j$, with their orders $m_j$, are exactly
\begin{equation}\label{eq:v5-test-list}
\begin{array}{c|c}
 f(U_0),\ldots,f(U_4),\ \sigma f(V_0),\ldots,\sigma f(V_4)&1\quad(10)\\
 L_i-L'_i\ (1\le i\le4)&16,52,52,52\\
 L_1-\beta&16\\
 F(L_2,L_3,L_4)&52,
\end{array}
\end{equation}
where
\begin{equation}\label{eq:v5-lucas-polynomial}
 F(U,V,W)=V^2-U(U-2)VW+(U-2)^2W^2-1.
\end{equation}
There is no rational proxy for an integer point label.

For $j=1,\ldots,4$ write $x(Q_j)=x_{j,0}+x_{j,1}\sqrt2$ and set
\[
 \mathcal A_0=1,\quad
 \mathcal A_{2j-1}=x_{j,0}^2,\quad
 \mathcal A_{2j}=(1+x_{j,1}^2)^2.
\]
These are positive nonzero rational squares. In fact $x(Q_j)$ is
positive at both real embeddings by Lemma~\ref{lem:v5-point-domains},
so $x_{j,0}>0$. Retain one rational coordinate $\lambda$ and put
\[
 \Theta=\{0,-1/2,-1,1/2399,-32/407\},\qquad
 F_5(T)=\prod_{\theta\in\Theta}(T-\theta),\qquad
 t_r=\mathcal A_r\lambda,\quad\delta_r=t_r-1.
\]
For every $0\le r\le8$, require $t_r\notin\Theta\cup\{1\}$ and define
\begin{align}
 R_r&=\mathcal A_r^{-5}F_5(t_r)/\delta_r^6,
 &\widetilde D_{r,j}&=D_j/\delta_r^{m_j+1},\label{eq:v5-detector-data}\\
 W_r&=\prod_{j=1}^{16}
 \left(1+R_r(\widetilde D_{r,j}/R_r^{m_j})^2\right),
 &T_r&=\Tr_{K/\Q}(W_r^{-1}),\qquad B_r=R_r/T_r^2.
 \label{eq:v5-detector-product}
\end{align}
Require $W_r,T_r\ne0$. For $B=B_r$ retain four rational witnesses
$e_r,s_r,p_r,q_r$ and impose
\begin{equation}\label{eq:v5-zero-pole}
 \begin{aligned}
 C_r&=17B_r,\qquad h_r=e_r(C_r-B_r^2),\\
 p_r^2&=C_r(1-e_r)(1+3e_r),\\
 C_rq_r^2&=(s_r^2-C_rh_r)(s_r^2+3C_rh_r).
 \end{aligned}
\end{equation}
The first two expressions are arithmetic definitions. The last two
are the preparation-zero specialization of the original pole predicate.
There is no nonzero guard on $e_r,s_r,p_r,q_r$, or $h_r$.

The scalar block shares the point $H$. Define
\[
 r_\beta=\beta(\beta-1)(\beta-2401)(32\beta+343)(2\beta-1).
\]
It is nonzero for rational $X$: the only positive possible root not
excluded by $\beta\ge2$ would require $X^2=2399$.
Let $Y_1,\ldots,Y_6$ be, in order, the two rational components of
$f(H),f(S),z(S)/z(H)$, and set
\[
 (e_1,\ldots,e_6)=(1,1,1,1,32,32),\qquad
 y_i=Y_i/r_\beta^{e_i},\qquad N_r(u,v)=u^2+rv^2.
\]
The scalar norm tree has no padding leaf. It is
\begin{equation}\label{eq:v5-scalar-tree}
 \begin{aligned}
 n_0&=N_{r_\beta}(1,y_1),\quad n_1=N_{r_\beta}(y_2,y_3),\quad
 n_3=N_{r_\beta}(y_5,y_6),\\
 W_\beta&=N_{r_\beta}\bigl(N_{r_\beta}(n_0,n_1),
                          N_{r_\beta}(y_4,n_3)\bigr).
 \end{aligned}
\end{equation}
Use eight Boolean coordinates $\epsilon_-,\epsilon_p$, for
$p\in\{2,3,5,7,11,13,17\}$, impose $\epsilon^2-\epsilon=0$, and set
\[
 c=(1-2\epsilon_-)\prod_p(1+(p-1)\epsilon_p),\qquad
 b_\beta=c r_\beta W_\beta^2,\qquad C_\beta=17b_\beta.
\]
Require $b_\beta\ne0$. Retain five rational witnesses
$t_\beta,e_\beta,s_\beta,p_\beta,q_\beta$ and impose
\begin{align}
 (1+b_\beta t_\beta^2)U_\beta
    &=b_\beta(1-b_\beta t_\beta^2),
 &h_\beta&=e_\beta(C_\beta-U_\beta^2),\label{eq:v5-scalar-preparation}\\
 p_\beta^2&=C_\beta(1-e_\beta)(1+3e_\beta),
 &C_\beta q_\beta^2&=(s_\beta^2-C_\beta h_\beta)
                     (s_\beta^2+3C_\beta h_\beta).
 \label{eq:v5-scalar-pole}
\end{align}
Here $U_\beta,h_\beta$ are arithmetic values. The preparation equation
forces $1+b_\beta t_\beta^2\ne0$: otherwise $b_\beta t_\beta^2=-1$
but its right side is $2b_\beta\ne0$.

Finally fix the integers
\begin{equation}\label{eq:v5-fixed-constant}
 C=3^{2^{26}},\qquad \mathsf S=C^2=3^{2^{27}},\qquad
 \kappa=\mathsf S^2=3^{2^{28}}.
\end{equation}
The symbol \texttt{S} in the coefficient files means $\mathsf S$;
it is not the point $S$ and is not a quantified coordinate.
Write $f(S)=f_0+f_1\sqrt2$, retain three private rational witnesses
$u,v,w$, and impose
\begin{equation}\label{eq:v5-physical-anchor}
 u^2+6v^2-2w^2-3(f_0^2+f_0f_1-f_1^2)^2+2\mathsf S^2=0.
\end{equation}
This completes the rational source. Every rational expression has the
displayed guard or a domain implication already proved. Equations
\eqref{eq:v5-zero-pole} and \eqref{eq:v5-scalar-pole} retain their
degenerate witness values.

\begin{lemma}[The three-witness anchor]\label{lem:v5-three-witness-anchor}
Equation~\eqref{eq:v5-physical-anchor} is equivalent to
$f_0/C,f_1/C\in\Z_{(2)}\cap\Z_{(3)}$. For the nonzero integer label
$a$ of $S=[a]P_*$ it implies
\[
 0\ne a\in2\,3^{2^{25}-1}\Z,\qquad
 |a|>(21H_*)^{13/2}>(20H_*)^4.
\]
Its permitted labels contain arbitrarily divisible positive integers.
\end{lemma}
\begin{proof}
Put $R=\Z_{(2)}\cap\Z_{(3)}$. The quaternion algebra $(2,3)_\Q$
ramifies exactly at 2 and 3. The quadratic embedding criterion
\cite[Proposition 14.6.7]{voight2021} shows, for rational $t\ne0$,
that $1-2/(3t^2)$ is represented by $2x^2+3y^2-6z^2$ precisely
when it is nonsquare in both $\Q_2$ and $\Q_3$. This last condition
is exactly $t\in R$. Indeed at 2 a positive valuation of $t$ gives
odd valuation of the expression, valuation zero gives residue 3
modulo 8, and negative valuation gives a square in $1+8\Z_2$.
At 3, nonnegative valuation of $t$ gives odd negative valuation,
while negative valuation gives a square in $1+3\Z_3$.
The expression is nonzero because $t^2=2/3$ has no rational solution.
Multiplying by $3t^2$ gives
\[
 t\in R\quad\Longleftrightarrow\quad
 \exists u,v,w\in\Q:\ u^2+6v^2-2w^2-3t^2+2=0.
\]
The case $t=0$ has witnesses $(0,0,1)$. A rational-square value of
$1-2/(3t^2)$ is excluded by the division-algebra criterion as well:
a nonzero trace-zero quaternion cannot square to a nonzero rational
square in a division algebra.

The form $A^2+AB-B^2$ has no nontrivial zero over $\F_2$ or $\F_3$.
Scaling a pair therefore proves
$v_p(a^2+ab-b^2)=2\min(v_p(a),v_p(b))$ for $p=2,3$.
Apply the preceding three-witness formula to
$(f_0^2+f_0f_1-f_1^2)/C^2$ and multiply its witnesses by $C^2$.
This gives exactly \eqref{eq:v5-physical-anchor}.

The local calculation in Section~\ref{subsec:o8-anchor-subgroup}
applies with exponent $2^{26}$: the reduction of $P_*$ at 3 has
order two and $v_3(z(2P_*))=1$. For even nonzero $a$ it gives
$v_3(f([a]P_*))=2+2v_3(a)$. At 2 the rational components of
$f([a]P_*)$ are integral. The anchor is consequently equivalent to
$0\ne a\in2\,3^{2^{25}-1}\Z$. Its minimum possible absolute value
has binary logarithm greater than
\[
 1+\tfrac32(2^{25}-1)=50331647.5,
\]
whereas
$\log_2((21H_*)^{13/2})<\tfrac{13}{2}(5+4783227)=31091008$.
The last comparison with $(20H_*)^4$ follows from
$21^{13}>20^8$ and $H_*\ge1$. The subgroup description proves the
last assertion.
\end{proof}

\subsection{Local tests and the height implication}
\label{subsec:v5-soundness}

\begin{lemma}[One place for every test]\label{lem:v5-one-place}
Let $R\in\Q^*$, $D_j\in K$, and positive integers $m_j$ define
\[
 W=\prod_j(1+R(D_j/R^{m_j})^2),\quad
 T=\Tr(W^{-1}),\quad B=R/T^2,
\]
with $W,T\ne0$ and a preparation-zero pole certificate for $B$.
If $p>76176$ and $v_p(R)$ is odd, then $v_p(R)>0$ and some
$\mathfrak p\mid p$ has $v_{\mathfrak p}(D_j)\ge m_j$ for every $j$.
\end{lemma}
\begin{proof}
Set $\alpha=v_p(R)$. The nonconstant term in each nontrivial factor
has odd valuation $\alpha+2(v(D_j)-m_j\alpha)$; unramifiedness
makes these valuations integral. The factor valuation is its minimum
with zero. Thus $\lambda_{\mathfrak p}=v_{\mathfrak p}(W)\le0$.
Put $\lambda=\max_{\mathfrak p\mid p}\lambda_{\mathfrak p}$.
Integrality of trace gives $v_p(T)\ge-\lambda$. Hence
\[
 v_p(B)\le\alpha+2\lambda.
\]
The left side is odd and, by Lemma~\ref{hyp:o8-pole}, not negative.
It is positive, so $\alpha>0$ and $\lambda>-\alpha/2$.
At a place attaining $\lambda$, each factor has valuation at least
$\lambda$, because all factor valuations are nonpositive. For
nonzero $D_j$ this implies
\[
 v(D_j)>(m_j-3/4)\alpha\ge m_j-3/4.
\]
Integrality gives $v(D_j)\ge m_j$; a zero $D_j$ satisfies it as well.
\end{proof}

\begin{lemma}[The two-anchor comparison]\label{lem:v5-two-anchors}
Write the integer labels of $(A,S,T,Z)$ as $(n,a,t,b)$.
If the first fourteen tests in \eqref{eq:v5-test-list} have their
specified orders at a place $\mathfrak p\mid p$, $p>76176$, then
\[
 L_i\in\mathcal O_{K_{\mathfrak p}},\qquad
 L_i\equiv n_i\pmod{p^{k_i}},\qquad
 (n_1,n_2,n_3,n_4)=(n,a,t,b).
\]
This holds at split and inert primes.
\end{lemma}
\begin{proof}
Write $H=[h]P_*$ and $U_i=[u_i]P_*$. The order-one tests put
$H,U_i$ in the formal group at $\mathfrak p$ and $V_0,V_i$ in the
formal group at the conjugate place. Positive valuation of $f=1/x$
gives formal membership on this good model. The guarded points give
$h\ne0,4$. Their full logarithm ratios, after conjugating the second,
are
\[
 r_i=u_i/h,\qquad r'_i=(u_i-4n_i)/(h-4).
\]
These identities follow from the homomorphism property on equal
formal multiples, and do not require $P_*$ itself to be formal.
At least one of $h,h-4$ is a $p$-unit. Its ratio is integral, and
Lemma~\ref{lem:v5-unramified-log} gives the same valuation to its
truncated ratio. The cross test makes the other truncated ratio
integral too. Both relative errors are therefore absolute errors,
and $L_i\equiv r_i\equiv r'_i$ modulo $p^{k_i}$. Multiplication
by $h,h-4$ and subtraction give $4L_i\equiv4n_i$. Divide by four.
At an inert prime conjugation preserves the single local valuation,
so the same argument applies.
\end{proof}

\begin{lemma}[Height transfer]\label{lem:v5-height-transfer}
Let $M_r$ be the odd-contact radical of $t_r$ for the five contacts
$\Theta$, omitting primes at most 76176. If $M_r^k$ divides a
nonzero integer $N$ for every $r$, then
\[
 h(x(Q_i))<7H_*|N|^{24/k}\qquad(1\le i\le4).
\]
At every prime counted in $M_r$, the source imposes all sixteen
field tests with their specified orders at one place above that prime.
\end{lemma}
\begin{proof}
The fractional linear map $\phi(t)=2+1/t$ carries $\Theta$ to
$\infty,0,1,2401,-343/32$. Its integral matrix has determinant $-1$.
All primes dividing the relevant fixed coefficients and contact
differences, including differences from 1, lie in
$\{2,3,5,7,11,37,109,439,2399\}$. Thus outside the cutoff
$M_r=M_5(\phi(t_r))$, and
\[
 h(t_r)\le h(\phi(t_r))+\log3
       \le(H_*+\log3)M_r^{24}.
\]
At a prime of $M_r$, $\delta_r$ is a unit. Exactly one contact
factor has odd positive valuation, while $\mathcal A_r^{-5}$ is
a rational square. Therefore $v_p(R_r)$ is odd.
Lemma~\ref{lem:v5-one-place} applied to $\widetilde D_{r,j}$ gives
the asserted orders for $D_j$.

Set $T_0=(H_*+\log3)|N|^{24/k}$. The case $r=0$ gives
$h(\lambda)\le T_0$, and the other cases give
$h(\mathcal A_r\lambda)\le T_0$. Hence
$h(\mathcal A_r)\le2T_0$. For the first component,
$h(x_{i,0})\le T_0$. For $x_{i,1}=a/b$ in lowest terms,
\[
 h(1+x_{i,1}^2)=\log(a^2+b^2)\ge2h(x_{i,1}),
\]
so $h(x_{i,1})\le T_0/2$. Finally
\[
 h(x(Q_i))\le h(x_{i,0})+h(x_{i,1})+\tfrac32\log2
 <7H_*|N|^{24/k}.
\]
\end{proof}

\begin{lemma}[The two reduced group cardinalities]
\label{lem:v5-two-reductions}
Let $p>76176$ be prime and let $\pi\in\Z[i]$ have norm $p$, with
$\pi\equiv9\pmod{16}$ and $\pi\equiv1+3i\pmod7$.
The two reductions of $E/K$ at the places above $p$ have group
cardinalities $N(\pi-1)$ and $N(\pi+1)$, in some order.
\end{lemma}
\begin{proof}
The congruences give $p\equiv1\pmod8$ and $p\equiv3\pmod7$.
Thus $p$ splits in $K$ and $(7/p)=-1$. The reductions of
$d=3-\sqrt2$ have product 7 and opposite square classes. The two
curves $y^2=x^3-d^2x$ are consequently opposite quadratic twists
of $y^2=x^3-x$ and have opposite traces.

The Frobenius calculation already given in
Section~\ref{subsec:o8-pattern-proof} applies: the order-four
automorphism is defined over $\F_p$, its rational centralizer is
$\Q(i)$, and integrality puts Frobenius in $\Z[i]$.
Write it as $a+bi$, so $a^2+b^2=p$ and the two cardinalities are
$p+1-2a$ and $p+1+2a$. Full rational two-torsion makes both
cardinalities divisible by four. Since $p\equiv1\pmod4$, this
forces $a$ odd and $b$ even. The real part of $\pi$ is odd and its
imaginary part is even. Unique factorization in $\Z[i]$ therefore
makes $a+bi$ one of $\pm\pi,\pm\overline\pi$; the associates
with factor $i$ have the wrong parity. Hence the unordered
cardinalities are exactly $N(\pi-1),N(\pi+1)$.
\end{proof}

\begin{theorem}[Arithmetic interpretation of the source]
\label{thm:v5-source-projection}
With the inherited arithmetic results and
Theorem~\ref{thm:v5-simultaneous-primes}, the rational source of
Section~\ref{subsec:v5-complete-source} has input projection exactly $\Z$.
\end{theorem}
\begin{proof}[Soundness]
Let $\mathcal B=\max(|n|,|a|,|t|,|b|)$ and
$\Delta=F(a,t,b)\in\Z$. Polynomial substitution preserves
congruences between integral local elements. Lemmas
\ref{lem:v5-two-anchors} and \ref{lem:v5-height-transfer} therefore
give $M_r^{52}\mid\Delta$ whenever $\Delta\ne0$.
The index-height estimate \eqref{eq:o8-indexheight} yields
\[
 \mathcal B^2-6<7H_*|\Delta|^{24/52},\qquad
 |\Delta|<5\mathcal B^4,
\]
the latter because $\mathcal B\ge|a|>7$. Consequently
\[
 \mathcal B^{2/13}<(7\cdot5^{6/13}+6)H_*<21H_*.
\]
For the last strict inequality, $5^6 7^{13}<15^{13}$.
This contradicts Lemma~\ref{lem:v5-three-witness-anchor}. Thus
$\Delta=0$, and Lemma~\ref{lem:o8-lucas}, with $b\ne0$, gives
$\log\mathcal B>|a|-3$.

Lemma~\ref{v5:scalar32-bridge} applies to the exact tree
\eqref{eq:v5-scalar-tree} and gives
\[
 h(\beta)\le H_*|a|^{3/4}<|a|/20.
\]
If $\beta\ne n$, put $N=|\operatorname{num}(\beta-n)|$ in lowest
terms. The order-16 test gives $M_r^{16}\mid N$. Therefore
\[
 \mathcal B^2<7H_*N^{3/2}+6,\qquad
 N\le2e^{h(\beta)}\mathcal B,
\]
and $\log\mathcal B<2\log(32H_*)+3h(\beta)$.
The anchor also gives $\log(32H_*)<|a|/20$, so the right side
is less than $|a|/4<|a|-3$, a contradiction. Hence
$\beta=n\in\Z$. Since $X\in\Q$ and $X^2=n-2\in\Z$,
the reduced denominator of $X$ is one.
\end{proof}

\begin{proof}[Completeness]
Fix $X\in\Z$, set $n=X^2+2$, and choose the Boolean correction
$c$ so that $cr_\beta>0$ is square in $\Q_2,\Q_3,\Q_{17}$,
as in the finite squareclass computation accompanying
Lemma~\ref{hyp:o8-pole}. Let $\mathcal S$ be the prime support of
$cr_\beta$. The scalar completeness construction in
Lemma~\ref{v5:scalar32-bridge} first chooses $h_0>0$ divisible by
64 with $[h_0]P_*$ sufficiently deep at these primes. It then
chooses a positive label $a$ sufficiently divisible relative to
$h_0$ and satisfying the anchor. The normalized scalar leaves are
integral at primes dividing $r_\beta$ and divisible at primes
dividing $c$ but not $r_\beta$. Choose positive $t,b$ with
$F(a,t,b)=0$ by Lemma~\ref{lem:o8-lucas}.
This fixes the four root points and all scales $\mathcal A_r$.

Temporarily put $H=[h]P_*$ and $U_i=[hn_i]P_*$.
All comparison tests are rational functions of $[h]P_*$ with
nonzero denominators at $h=h_0$. Their poles lie above finite
sets under fixed nonzero multiplication and translation maps
on the two real factors, as required by the real progression lemma.
Continuity preserves the scalar
conditions at every place over $\mathcal S$ and bounds below the
valuations of all sixteen comparison tests at the places over
$2,3,17$, throughout a progression
\[
 h\equiv h_0\pmod L,\qquad64\mid L.
\]
One obtains $L$ by putting sufficiently small open formal subgroups
inside the finitely many translated neighborhoods and taking a
common multiple of their indices. Lemma~\ref{v5:real-progression-bridge}
gives, independently of future progression moduli, real bounds for
the tests attained by arbitrarily large positive $h$ in each
nonconstant progression. The direct polynomial $F$ introduces
no denominator and has the same boundedness property.

Choose a positive integer $D_0$ such that every
$D_0^2\theta/\mathcal A_r$ is integral, and put
\[
 v=D_0^2w^2,\qquad\lambda=u/v.
\]
The contact positions are the integer forms
$u-(D_0^2\theta/\mathcal A_r)w^2$. Delete repeated forms, but retain
their incidences in the nine detectors. The five forms within one
detector are distinct, and the zero contact is the common form $u$.
Apply Theorem~\ref{thm:v5-simultaneous-primes} to these forms.
Prescribe arbitrarily high divisibility of $w$, a sufficiently
large lower bound for $u/w^2$, and primes larger than 76176 and
all prime factors of the fixed scales and contact differences.
In its odd squarefree avoidance modulus include every odd prime
factor of $3\cdot17L$ except 7.
Writing this modulus as $M$, the prescribed Gaussian classes are
\[
 \pi_p\equiv9\pmod{16},\qquad
 \pi_p\equiv1+3i\pmod7,\qquad
 \pi_p\equiv i\pmod M.
\]

Let the resulting Gaussian generators be $\pi_p$, and write
$d_-(p)=N(\pi_p-1)$ and $d_+(p)=N(\pi_p+1)$ for the cardinalities
of the two reduced elliptic groups, orienting the places by
Lemma~\ref{lem:v5-two-reductions}. These cardinalities annihilate
the reductions of $P_*$; their exact point orders are not needed.
They satisfy
\[
 v_2(d_-)=6,\quad v_2(d_+)=2,\quad
 \gcd(d_-(p),d_+(q))=4
\]
for every pair of selected primes. No odd prime factor of $L$
divides either order. For 7 this follows separately from
$d_+\equiv6$ and $d_-\equiv2\pmod7$.
Set $M_-=\operatorname{lcm}_p d_-(p)$ and
$M_+=\operatorname{lcm}_p d_+(p)$. The congruences
\[
 h\equiv0\pmod{M_-},\quad h\equiv4\pmod{M_+},\quad
 h\equiv h_0\pmod L
\]
are compatible: their pairwise gcds are $4,64,4$, respectively.
Choose a sufficiently large positive $h$ in this progression
attaining the real bounds. It avoids $0,4,-4$, so all point and
subtraction guards hold. At a selected prime, $H,U_i$ are formal
at the place of order $d_-$, while $V_0,V_i$ are formal at its
conjugate place. Their exact ratios are $n_i$. Lemma
\ref{lem:v5-unramified-log} gives all sixteen required orders,
including the order-52 polynomial test because $F(a,t,b)=0$.

We check that the nine preparation-zero poles can be filled.
For each detector,
\[
 R_r=\frac{\prod_{\theta\in\Theta}
                   (u-(\theta/\mathcal A_r)v)}
                 {v^5\delta_r^6}.
\]
The denominator is a rational square. At the five distinct contact
primes of this detector $R_r$ has valuation one, and elsewhere
its valuation is even. The primes do not divide $v$: otherwise
two distinct forms in the same detector would give the same
prime value, impossible since $w>0$ and their coefficients are
distinct. The lower bound for the primes excludes the fixed scale
and difference denominators. At either place over a
contact prime a factor of $W_r$ has valuation zero or negative,
and its inverse has residue one or zero, respectively. At the
successful place every inverse factor has residue one.
Consequently $T_r$ has residue one or two and is a unit.
It follows that $B_r$ has positive odd valuation at its contact
primes and even valuation elsewhere. Its reduced denominator is
a square.

The prescribed Gaussian classes make every contact prime square
in $\Q_2,\Q_3,\Q_{17}$. The last product formula therefore makes
$B_r$ square in each of those fields. The uniform real bounds
and $\lambda\longrightarrow+\infty$ give
\[
 R_r\sim\mathcal A_r^{-6}\lambda^{-1},\quad
 W_r=1+O(\lambda^{-3}),\quad T_r=2+O(\lambda^{-3}),\quad
 B_r\sim(4\mathcal A_r^6\lambda)^{-1}.
\]
In particular $0<B_r<17$. At a small prime $p=2,3,17$, put
$\nu_p=v_p(v)$ and $\eta_{p,r}=v_p(\mathcal A_r)$.
The common form $u$ is a selected prime greater than these small
primes, so $v_p(u)=0$ and $v_p(\lambda)=-\nu_p$.
The fixed neighborhoods give $v(D_j)\ge-c_p$. For sufficiently
large prescribed $\nu_p$ one has
\[
 v_p(R_r)=\nu_p-6\eta_{p,r},\qquad
 v(\widetilde D_{r,j})\ge-c_p+(m_j+1)(\nu_p-\eta_{p,r}).
\]
The nonconstant factor term has valuation at least
$3\nu_p+(10m_j-8)\eta_{p,r}-2c_p$, tending to infinity.
Thus $T_r$ is as close to 2 as needed, including at the ramified
prime 2, and
\[
 v_p(B_r)=\nu_p-6\eta_{p,r}-2v_p(2).
\]
Prescribing the divisibility of $w$ before applying the prime theorem
therefore ensures $v_2(B_r)\ge3$, $v_3(B_r)\ge1$,
$v_{17}(B_r)\ge2$. Lemma~\ref{v5:preparation-zero-bridge} supplies
the four witnesses in every coordinate pole.

Finally the progression preserved all scalar leaf conditions.
The constant leaf on the all-left path makes $W_\beta$ a unit at
every prime dividing $r_\beta$; the divisibility of the other
leaves makes it a unit at primes dividing $c$ but not $r_\beta$.
Since $|r_\beta|>1$, it is nonzero. Therefore
$b_\beta=cr_\beta W_\beta^2$ is positive, has square reduced
denominator, and is locally square at $2,3,17$.
Lemma~\ref{hyp:o8-pole} supplies the scalar preparation and pole
witnesses for this final $H$. The anchor already holds for $S$.
All witnesses have now been supplied without changing the root points.
\end{proof}

% Replace the former exact-compilation subsection with this fragment.
% Retain the preceding shared-point arithmetic source and its hypotheses.
% No claim here removes or proves its simultaneous-prime input.
\subsection{The quadratic compilation in 962 variables}
\label{subsec:v6-962-compilation}

We compile the shared-point source above using different point charts,
pole coordinates, and arithmetic circuits. Every change has explicit
rational witness maps. The arithmetic source, including its simultaneous
prime input, is unchanged. Throughout this subsection
\[
 K=\Q(\sqrt2),\qquad D=(3-\sqrt2)^2=11-6\sqrt2,
 \qquad \mathsf S=3^{2^{27}}.
\]
An equation in $K$ means its two rational component equations. The fixed
integer $\mathsf S$ is a coefficient, not a coordinate. An assignment of
a quadratic expression to one rational coordinate adds one quadratic
equation. Affine expressions require no coordinate.

\begin{theorem}[The 962-variable quartic]
\label{thm:v6-962-polynomial}
The complete table specified below defines
\[
 P_{962}\in\Z[X_0,\ldots,X_{961}]
\]
with degree four, exactly $55409$ nonzero coordinate monomials, integer
coefficient content one, and coefficient height
\[
 H(P_{962})=64\mathsf S^2=64\cdot3^{2^{28}}<2^{429496736}.
\]
It is the sum of squares of $919$ quadratic residuals and has the same
rational input projection as the shared-point source of
Theorem~\ref{thm:v5-source-projection}. Consequently, under that theorem's
arithmetic hypotheses, for every $x\in\Q$,
\[
 x\in\Z\quad\Longleftrightarrow\quad
 \exists X_1,\ldots,X_{961}\in\Q,\quad
 P_{962}(x,X_1,\ldots,X_{961})=0.
\]
The $962$ variables include the input and $961$ existential witnesses.
\end{theorem}

\paragraph{A quadratic chart for the nine points.}
For each point introduce $t,h\in K$ and set
\begin{equation}\label{eq:v6-point-chart}
 u=t^2,\qquad X=h^2,\qquad
 u(X-u/4)=D,\qquad Y=h(u/2-X).
\end{equation}
The three assignments use six rational coordinates. The displayed
constraint contributes two quadratic residuals. The four rational
components of $t,h$ replace the four halfpoint components.

\begin{lemma}[Exact point-chart maps]\label{lem:v6-point-chart}
The chart \eqref{eq:v6-point-chart} describes precisely the finite doubles
used in the shared-point source, with the same rational witnesses up to
the following invertible changes. It preserves all required nonzero
point-coordinate domains. The rational component $X_0$ of $X$ is positive.
\end{lemma}
\begin{proof}
For an old halfpoint $(a,b)$ with $b^2=a^3-Da$ and its tangent equation,
both $a,b$ are nonzero. Indeed $b=0$ would give $a=0$ or $a^2=D$,
whereas the tangent numerator $3a^2-D$ would be $-D$ or $2D$. Then
$a=0$ is excluded by the curve equation. Define
\[
 t=b/a,\qquad h=(a^2+D)/(2b)=a/t-t/2.
\]
The curve equation gives $t^2=a-D/a$. Substitution gives
$t^2(h^2-t^2/4)=D$. The doubling formulas give exactly $X=h^2$
and $Y=h(t^2/2-h^2)$.

Conversely, a chart solution has $t\ne0$ because $D\ne0$. Set
\[
 a=th+t^2/2,\qquad b=at,\qquad m=h+t.
\]
If $a=0$, then $h=-t/2$, contradicting the chart equation. Thus $b\ne0$.
Direct expansion gives
\[
 a^2-t^2a=D,\quad b^2=a^3-Da,\quad
 2bm=3a^2-D,\quad m^2-2a=h^2,
\]
and $m(a-h^2)-b=h(t^2/2-h^2)$. The substitutions are inverse, so they
recover exactly the old finite double. Lemma~\ref{lem:v5-point-domains}
therefore supplies its nonzero abscissa and ordinate, and the same remains
true for the unchanged guarded differences and subsequent doublings.

Both real embeddings of $D$ are positive. Since $t\ne0$, the chart
equation implies $h^2=t^2/4+D/t^2>0$ under both embeddings. Hence
$X_0$, the average of these two positive values, is positive.
\end{proof}

The five affine subtraction formulas keep their original inverses of
the differences of abscissas. Neither a subtraction guard nor a point
occurrence is deleted. Relative to the direct halfpoint doubling circuit,
the chart saves two rational arithmetic coordinates for each of the nine
points.

\paragraph{The logarithm polynomial.}
Retain the source's coefficients
\[
 a_0=1,\quad a_1=-2/5,\quad
 a_j=a_{j-1}(4j-5)/(4j+1).
\]
For $P(w)=\sum_{j=0}^{12}a_jw^j$ write
\begin{equation}\label{eq:v6-log-square-completion}
 P(w)=a_{12}Q(w)^2+R(w),\qquad
 \deg Q=6,\quad Q\text{ monic},\quad\deg R\le5.
\end{equation}
This identity is constructed without any witness division. Starting
with the leading coefficient of $Q$ equal to one, match degrees eleven
down to six successively. At each step the new coefficient is determined
by division by $2a_{12}\ne0$. The remainder then has degree at most five.
After computing $w^2,w^3$, every term through $w^6$ is quadratic in
$1,w,w^2,w^3$. Assigning $Q$ and then $a_{12}Q^2+R$ therefore uses
four field assignments, rather than five. The eight order-$52$
logarithms save sixteen rational coordinates. The order-$16$ logarithms
keep their exact degree-three polynomials. No truncation order changes.

Usually the final logarithm is $\ell=zP(Df^2)$, with
$f=1/x$ and $z=-x/y$. For $U_0$ and $V_0$ only, omit the separate $z$
coordinates and impose the direct quotient equations
\begin{equation}\label{eq:v6-log-quotient}
 y\ell=-xP(Df^2).
\end{equation}
The polynomial output is an affine expression in assigned coordinates,
so these are quadratic component equations. The ordinate is nonzero by
Lemma~\ref{lem:v5-point-domains}; hence \eqref{eq:v6-log-quotient}
is exactly the old logarithm. The source's logarithm-nonvanishing result
and every logarithm-ratio domain remain unchanged. In the scalar block
replace $z(H)\rho=z(S)$ by
\begin{equation}\label{eq:v6-scalar-ratio}
 x(H)\rho=-y(H)z(S).
\end{equation}
Both $x(H),y(H)$ are nonzero, so this has the same unique value
$\rho=z(S)/z(H)$.

\paragraph{Products, tests, and detector evaluation.}
For each of the three orders, compute the elementary symmetric
coefficients of $\prod_j(1+qD_j^2)$ by a balanced product tree.
At an internal node with coefficient vectors $a_i,b_j$, assign each
component of $\sum_{i+j=k}a_ib_j$ as one quadratic expression.
This is the ordinary coefficient-convolution identity. It preserves
every coefficient and every test, including its assigned order.

The degree-ten detector polynomial is evaluated in two blocks:
\begin{equation}\label{eq:v6-detector-block-evaluation}
 \sum_{j=0}^{10}c_jq^j
 =\sum_{j=0}^{4}c_jq^j
   +q^5\left(\sum_{j=0}^{5}c_{j+5}q^j\right).
\end{equation}
Compute $q^2,q^3,q^4,q^5$ using four scalar assignments, assign the
two components of the parenthesized block, and assign the two final
components directly. The low block is not assigned separately.
Every final summand is quadratic in the available coordinates. This
uses eight rational coordinates. The degree-two and degree-four
polynomials use their direct power sums, with their two component
sums assigned once. These formulas are identities over $K$.

The Lucas expression is regrouped as
\[
 t^2-a(a-2)tb+((a-2)b)^2-1
 =t^2-(at)C+C^2-1,\qquad C=(a-2)b.
\]
Thus one field multiplication is removed without changing the tested
polynomial. The fixed quintic is evaluated by its exact coefficients
using $t^2$; it remains
$F_5(t)=t(t+1/2)(t+1)(t-1/2399)(t+32/407)$.

\paragraph{Weight quotients and their domains.}
Fix a detector and write $A$ for its scale, $\delta=A\lambda-1$,
and $e$ for its weight. The scales are
\[
 1,\quad x(Q)_0^2,\quad (1+x(Q)_1^2)^2
 \qquad(Q=A,S,T,Z).
\]
They are all positive by Lemma~\ref{lem:v6-point-chart}. Retain the
inverse $\delta i=1$, and replace the separate fifth-power product
and inverse of $F_5$ by
\begin{equation}\label{eq:v6-weight-domain}
 F_5(A\lambda)e=(A\delta)^4(A\delta).
\end{equation}
Here $(A\delta)^4$ denotes the output of its retained power-assignment
gates; the displayed product is quadratic in the available coordinates.
The right side is nonzero, so this equation forces both $F_5(A\lambda)$
and $e$ nonzero. In particular the removed inverse of $F_5$ is recovered
uniquely. Define the low-order weight by
\[
 \delta q_1=e i^2,
\]
where $i^2$ is likewise a retained square-assignment output. This gives
the unchanged $q_1=e/\delta^3$. Retain the repeated-squaring
values $e^6,e^{24},e^{96}$ and set
\[
 c=q_1e^6,\qquad q_{16}=ce^{24},\qquad q_{52}=ce^{96}.
\]
These are the old $q_1e^{30}$ and $q_1e^{102}$, with one shared product.

If the detector product is $W=w_0+w_1\sqrt2$, retain $w_0j=1$ and use
\begin{equation}\label{eq:v6-norm-final-weight}
 w_0n=(w_0^2-2w_1^2)/2,\qquad ev=in,\qquad B=nv.
\end{equation}
The trace guard is essential and is retained. It gives $W\ne0$ in $K$,
so $\operatorname{Norm}(W)\ne0$ and
$n=\operatorname{Norm}(W)/(2w_0)\ne0$. Equations
\eqref{eq:v6-weight-domain}--\eqref{eq:v6-norm-final-weight} therefore
give
\[
 B=\frac{i}{e}\left(\frac{\operatorname{Norm}(W)}{2w_0}\right)^2\ne0,
\]
the original detector input. The reverse map restores $R=i/e$, its
old inverse intermediates, and $n^2$ by their unique rational values.
No inverse is inferred from an unguarded zero-over-zero quotient.

\paragraph{The scalar norm and finite corrections.}
The scalar contact polynomial $r$ and its inverse $ri=1$ remain.
Writing $f=f(H)_0$, $u=f(H)_1$, and $g=f(S)_0$, retain the leaf
$\ell=ui$ and replace two norm-tree
nodes by
\[
 rn_0=r+f^2,\qquad rn_1=u\ell+g^2.
\]
with the existing leaf $ui$ retained. Since $r\ne0$, these equal
$n_0=1+r(fi)^2$ and $n_1=(ui)^2+r(gi)^2$. All other nodes and all
orders in the scalar norm tree are unchanged.

Replace the eight Boolean correction choices and their seven product
coordinates by $c_0,\ldots,c_7\in\Q$ satisfying
\begin{equation}\label{eq:v6-correction-chain}
 c_0^2=1,\qquad
 (c_j-c_{j-1})(c_j-p_jc_{j-1})=0\quad(1\le j\le7),
\end{equation}
where $(p_1,\ldots,p_7)=(2,3,5,7,11,13,17)$. Induction shows
$c_j\ne0$ and $c_j/c_{j-1}\in\{1,p_j\}$. Thus $c_7$ is exactly an
allowed signed squarefree correction. Conversely every original
choice determines such a chain. Recover the original Booleans by
\[
 \epsilon_-=(1-c_0)/2,\qquad
 \epsilon_j=\frac{c_j/c_{j-1}-1}{p_j-1}.
\]
All denominators are nonzero. The scalar preparation coordinate,
its defining quotient, and the explicit inverse of its pole input
are retained.

\paragraph{The ten pole blocks.}
Write $b$ for a pole input, $C=17b$, and $U$ for its prepared value.
For the nine coordinate poles $U=b$. In the scalar pole $U$ retains
the source's preparation formula and its denominator domain.
Every $C$ is nonzero: this was proved above for coordinate detectors,
and the scalar input has its retained inverse.

Replace the old witness $e_0$ by $z=Ce_0$. Introduce the five
intermediates
\[
 P=p^2,\quad F=U^2,\quad H=z(C-F),\quad V=s^2,\quad Q=q^2,
\]
and impose
\begin{equation}\label{eq:v6-scaled-poles}
 CP=(C-z)(C+3z),\qquad CQ=(V-H)(V+3H).
\end{equation}
Every equation is quadratic. Under $z=Ce_0$, the first residual is
$C$ times the old residual $p^2-C(1-e_0)(1+3e_0)$, and the second
is unchanged because $H=Ce_0(C-U^2)$. The inverse map $e_0=z/C$
is defined. This saves one arithmetic coordinate at each pole,
without changing the pole predicate or any preparation condition.

\paragraph{Counts and the final polynomial.}
The initial allocation is
\[
\begin{array}{l|r|r}
 &\text{coordinates}&\text{structural residuals}\\\hline
 \text{input}&1&0\\
 \text{nine point charts}&36&18\\
 \text{common slope}&1&0\\
 \text{nine coordinate poles}&36&18\\
 \text{scalar pole with preparation}&5&2\\
 \text{private anchor}&3&1\\
 \text{signed correction chain}&8&8\\\hline
 \text{total}&90&47
\end{array}
\]
The complete program has $904$ arithmetic coordinates and their
$904$ defining equations. Hence the raw counts are $994$ coordinates
and $951$ quadratic residuals.

The private anchor uses exactly the old affine change
\[
 u=\mathsf S+(U-3V+3W)/4,\quad
 v=\mathsf S/2+V/4,\quad w=3\mathsf S/2+W/4.
\]
Its determinant is $1/64$. If $T=f_0^2+f_0f_1-f_1^2$, sixteen times
the old anchor after substitution is
\[
 -48T^2+U^2-6UV+6UW+15V^2-18VW+7W^2+8\mathsf S U.
\]
No anchor condition or witness is removed.

Each certified compiler elimination solves a row $aY+q=0$ with
$a\in\Q^*$ and $q$ independent of $Y$, substitutes $Y=-q/a$
everywhere, and removes that row. It is accepted only if every
remaining residual is still quadratic. The reverse map reinserts
the unique displayed value of $Y$. There are exactly $32$ such
eliminations, no row additions, and no omission of an unused coordinate.
The input is retained. Consequently
\[
 90+904-32=962,\qquad47+904-32=919.
\]
Normalize the remaining rows to primitive integer polynomials
$q_1,\ldots,q_{919}$ and define
\begin{equation}\label{eq:v6-final-polynomial}
 P_{962}(X_0,\ldots,X_{961})=\sum_{\nu=1}^{919}q_\nu^2.
\end{equation}
The complete coefficient table is
\nolinkurl{quartics/quartic-962/polynomial.json}, with SHA-256
\begin{center}\small
\nolinkurl{6b798de7d00ac3c7d725ca0a69d39f1eac1ddb02fee2756cb6323770b952af32}.
\end{center}
An entry with coefficient $c$, exponent \texttt{sPower} $k$, and
coordinate list $[i_1,\ldots,i_d]$ means
$c\mathsf S^kX_{i_1}\cdots X_{i_d}$. Every entry is included and repeated
indices give powers. This defines an ordinary integer polynomial.

\begin{proof}[Proof of Theorem~\ref{thm:v6-962-polynomial}]
For a source solution apply the point-chart, correction-chain and
pole substitutions, and evaluate all new assignment and quotient
coordinates in dependency order. The identities and domain proofs
above show that every structural residual vanishes. Conversely apply
the inverse substitutions, restore every deleted quotient by its
proved nonzero denominator, and evaluate the source assignments.
Every source point, logarithm, test, detector input, scalar value and
anchor value is recovered. These maps fix $X_0$. The affine anchor
change, certified eliminations and nonzero row normalizations preserve
the same projection. Over $\Q$, \eqref{eq:v6-final-polynomial} vanishes
exactly when all $919$ residuals vanish. This proves the projection
statement, and Theorem~\ref{thm:v5-source-projection} gives its conditional
arithmetic interpretation.

Independent exact convolution gives $55409$ distinct monomials.
Of their coefficients, $55401$ are $\mathsf S$-independent integers,
seven are $c\mathsf S$ with $0<|c|\le768$, and one is
$64\mathsf S^2$. The largest absolute $\mathsf S$-independent coefficient
has bit length $1526$. Thus no coefficient vanishes on specialization
and $64\mathsf S^2$ is the exact height. It occurs at $X_{79}^2$,
where $X_{79}=U$. The coefficient of $X_0^4$ is one, proving degree
four and integer content one. Finally
$\log_2(64\mathsf S^2)<6+(8/5)2^{28}<429496736$ because $3^5<2^8$.
\end{proof}

The frozen package includes the builder, raw systems, all $904$ gate
definitions, the exact elimination and normalization certificate,
compact coordinate map, final residuals, complete coefficient table,
source-equivalence proofs and independent reviews. Its manifest has
SHA-256
\begin{center}\small
\nolinkurl{36fca72cb0eb126193ffa43d890ff09d5e29cbf1274561d2632ed1e227163787}.
\end{center}
The independent checker verifies every gate, all $894$ certificate
operations, the complete sum-of-squares expansion, support,
integrality, content and height. A fresh isolated reproduction gives
the same seven deterministic artifacts byte for byte. These finite
checks establish the stated algebraic facts; they do not prove any
inherited arithmetic input.
